%% Cover letter for submission to Environment Systems and Decisions (Springer).
%% Build: latexmk -pdf cover_letter.tex
%% Before sending: fill the editor's name if known, and add the conference
%% venue(s)/status of the prior work where indicated by the % TODO marks.
\documentclass[11pt]{article}
\usepackage[a4paper,margin=0.9in]{geometry}
\usepackage{parskip}
\usepackage{hyperref}
\hypersetup{hidelinks}
\setlength{\parskip}{4pt}

\begin{document}
\pagestyle{empty}

\noindent Alfa Yohannis\\
Universitas Pradita, Tangerang, Banten, Indonesia\\
\href{mailto:alfa.ryano@gmail.com}{alfa.ryano@gmail.com}

\vspace{0.6em}
\noindent \today

\vspace{0.6em}
\noindent The Editor-in-Chief\\
\emph{Environment Systems and Decisions}, Springer Nature

\vspace{0.6em}
\noindent Dear Editor, % TODO: replace with the editor's name if known

We are pleased to submit our manuscript, ``\textbf{Carbon Footprint and Carbon-Aware
Selection of Document-Oriented, Relational, and Hybrid Object-Relational Storage for
Image-Based Time-Series Workloads in Green IoT},'' for consideration as a research article
in \emph{Environment Systems and Decisions}. Smart buildings and monitoring systems now keep
streams of time-stamped images (for example, from cameras and sensors) in databases that run
around the clock, so how those images are stored is an environmental choice as much as a
technical one---yet its carbon cost has not been measured.

Using the same hardware and direct energy measurement throughout, we compared the carbon
footprint of three common storage approaches: a document database (MongoDB), a relational
database (PostgreSQL), and a hybrid that keeps only small descriptive records in the database
while placing each image in a separate file store. Testing image qualities from low resolution
up to very high resolution (6K), we found that no single option is best everywhere: the
document database is the lowest-carbon choice for ordinary image sizes, the hybrid becomes the
lowest-carbon choice for large images (roughly halving the emissions of storing images inside
the database), and the relational database sits consistently in between. At the largest size,
the document database cannot store the images at all, owing to a built-in limit on how large a
single record may be. We turn these measurements into a practical, carbon-aware guide for
choosing an approach, and we test how it holds up across countries with cleaner or dirtier
electricity, account for the manufacturing (``embodied'') carbon of the extra storage some
options consume, and project the yearly emissions of a realistic camera fleet.

The work suits \emph{Environment Systems and Decisions}: it treats an everyday infrastructure
choice as an environmental decision, combines direct measurement with both operational and
manufacturing carbon, considers how the answer depends on the local electricity grid, and
offers actionable guidance for sustainable building monitoring rather than a purely technical
benchmark. We also relate the findings to several UN Sustainable Development Goals.

In the interest of transparency, this manuscript consolidates and extends two of our earlier
conference papers % TODO: add venue(s)/status, e.g. "(presented at <venue>, <year>)"
that separately compared two of the three approaches. The present work measures the energy
directly (rather than estimating it from an assumed, fixed power draw), brings all three
approaches together under one common test, extends the range to the largest image size (where
the document-database limit appears), and adds the decision guide and the environmental
analyses that are its central contribution.

The manuscript is original, has not been published previously, and is not under consideration
elsewhere; all authors have read and approved it and declare no competing interests. To support
reproducibility, the benchmark code, configuration, and raw measurements will be released in a
public repository with a permanent identifier upon acceptance. Thank you for considering our
submission.

\vspace{0.6em}
\noindent Sincerely,

\vspace{1em}
\noindent Alfa Yohannis, on behalf of the authors\\
(Alfa Yohannis, Universitas Pradita; Alexander Waworuntu, Universitas Multimedia Nusantara)

\end{document}
